% sections/abstract.tex — Abstract
\begin{abstract}
The actor model has had a long influence on concurrent and
distributed systems design~\citep{hewitt1973actors,agha1986actors};
its classical formulations nevertheless assume CPU execution:
preemptive scheduling, heap-allocated mailboxes, and
operating-system-mediated supervision.
GPU programming has evolved in parallel along an orthogonal axis,
treating the device as a coprocessor that executes ephemeral kernels
launched by the host.  The two paradigms have remained largely
disjoint.  Recent hardware advances—NVIDIA Hopper's Thread Block
Clusters and Distributed Shared Memory~(DSMEM), cooperative kernel
launches with grid-wide synchronization, and stream-ordered
asynchronous allocation~\citep{nvidia2022hopper,nvidia2023cooperative}—make
a different model practical enough to formalize and evaluate
directly: \emph{persistent GPU actors} that live on the device
across the lifetime of the application, exchange messages through
device-side queues without host mediation, and recover from faults
under supervisor-driven restart policies.  Prior work has addressed
actor-to-accelerator integration, formal reasoning about GPU kernel
internals, and device-side task-based execution as three largely
separate threads; to our knowledge, no existing work combines
device-resident lifecycle, supervision, causal clocks, tenant
routing, and migration under a single formal account.

We present \system{}, a Rust-based runtime that realizes persistent
GPU actors, and---to our knowledge---the first formal model that
combines lifecycle, supervision, causal clocks, tenant routing, and
cross-GPU migration for device-resident actors.  Our contributions
are:
\textbf{(i)}~we \emph{define} a small-step operational semantics for
the actor lifecycle (\Dormant{}\,\(\to\)\,\Initializing{}\,\(\to\)\,%
\Active{}\,\(\to\)\,\Draining{}\,\(\to\)\,\Terminated{}, with
\Failed{} as a fault attractor) integrated with the GPU's SIMT
execution model and grid-wide cooperative synchronization;
\textbf{(ii)}~we define a kernel-to-kernel (K2K) messaging calculus
with explicit delivery semantics (no-loss accounting for accepted
envelopes, per-(sender,\,destination) FIFO, idempotency-gated
at-most-once observable delivery, hop-bounded routing) and an
explicit dead-letter discipline distinguishing overflow drops,
invalid-target drops, policy evictions, and idempotency rejections;
\textbf{(iii)}~we define a Hybrid Logical Clock (HLC) integration
that provides causal ordering across actors on a single device and
\emph{prove} that snapshot-restore preserves causality across
supervisor-driven restart;
\textbf{(iv)}~we define a tenant-isolation model in which K2K
sub-brokers are keyed on a two-tuple audit boundary
(\(\mathit{org\_id}, \mathit{engagement\_id}\)) and prove that the
supervisor's routing function admits no cross-tenant trace under
the model—a routing-level structural guarantee, not an end-to-end
confidentiality claim against hardware-capable adversaries;
\textbf{(v)}~we define a multi-GPU actor migration protocol with
phase-machine semantics; its atomic-swap, no-message-loss,
FIFO-across-cutover, and state-consistency properties are stated
and sketched analytically in \Cref{sec:migration} and checked
under bounded configurations in \TLA{}, and we state a
fast-path-selection property that NVLink-paired actors never fall
back to the host-staged path;
\textbf{(vi)}~we \emph{mechanize} the calculus of each preceding
contribution as a \TLA{} specification and discharge its invariants
via bounded TLC model checking across six modules (the checks are
sound up to the stated parameter bounds; we do not claim scale-free
proofs);
\textbf{(vii)}~we \emph{measure} the runtime on NVIDIA H100~NVL
silicon.  The host cost to submit a burst of 1{,}000 commands
through mapped memory is 0.18\,\(\mu\)s wall-clock, \emph{amortized}
over the burst through the CPU's write-combining buffer; the same
burst issued as 1{,}000 \texttt{cuLaunchKernel} calls costs
1{,}583.1\,\(\mu\)s, so the host pays 8{,}698\(\times\) less time
submitting the burst on this path.  (Per-command values and the
amortized-versus-one-shot distinction are spelled out in
Table~\ref{tab:injection}.)  K2K (kernel-to-kernel)
message latency is 72--83\,ns for payloads up to 1~KiB through the
intra-block lock-free queue; sustained throughput is 5.10\,M
messages/second with a coefficient of variation of 0.66\% over
four 60-second trials; cluster-level synchronization is
0.628\,\(\mu\)s (2.98\(\times\) faster than grid-wide
synchronization).  On a 2\(\times\) H100~NVL system the NVLink
P2P migration path reaches 257.78\,GB/s unidirectional at 16~MiB
(85.9\% of the 12-link NVLink4 bundle's 300\,GB/s unidirectional
peak, equivalently 43.0\% of the 600\,GB/s bidirectional peak)
and runs 8.65\(\times\) faster than the host-staged fallback at
the same size.  The command-submission
comparison measures the control path---mapped-memory writes
versus host-driven launches, amortized over a 1{,}000-command
burst---not a runtime-vs-runtime execution-speed claim.  The
\TLA{} specifications backing each formal claim are open; running
them under TLC caught a real spec-level bug before it reached the
paper (\Cref{sec:disc:tlc-bug}).

The model relates three design patterns observed in production
GPU systems that are not usually treated together: persistent
kernels for reduced launch overhead~\citep{gupta2012persistent,%
aila2009raytrace}, device-side work queues for
producer--consumer pipelines, and supervisor trees for fault
containment.  We suggest that \emph{persistent GPU actors} can be
read as a unifying abstraction for these patterns.  The paper
closes by sketching how the model extends to multi-tenant,
multi-engagement workloads in regulated domains (audit,
compliance, safety-critical reasoning) where a routing-level
cross-tenant isolation guarantee is itself an auditable
deliverable.
\end{abstract}
